\documentclass[11pt]{article}
\usepackage[T1]{fontenc}
\usepackage[margin=24mm]{geometry}
\usepackage{amsmath,amssymb,booktabs,hyperref}
\newcommand{\F}{\mathbb F_2}
\newcommand{\Z}{\mathbb Z}
\newcommand{\R}{\mathbb R}
\newcommand{\Sq}{\operatorname{Sq}}
\newcommand{\SSq}{\mathsf S}
\newcommand{\dd}{d_{s_1}}
\newcommand{\lift}[1]{\widetilde{#1}}
\title{Cochain conventions for the independent space-group calculation}
\author{AllFSPT calculation note}
\date{September 2026}
\begin{document}
\maketitle

\section{Physical layers and the equations being solved}
In three spatial dimensions the decorated state is
$(n_1,n_2,n_3,\widehat\nu_4)$, where $n_1$ counts integer $p+ip$ sheets,
$n_2$ is the Majorana-chain decoration, $n_3$ is the complex-fermion
decoration, and $\widehat\nu_4$ is an additive phase. Their coefficients are
\begin{equation}
 n_1\in C^1(G_b;\Z_{s_1}),\quad n_2\in C^2(G_b;\F),\quad
 n_3\in C^3(G_b;\F),\quad\widehat\nu_4\in C^4(G_b;\R/\Z_{s_1}).
\end{equation}
The extension $\omega_2$ and character $s_1$ are held fixed. The current
space-group calculation uses $\omega_2=0$ and the binary orientation character
defined by $(-1)^{s_1(g)}=\det(g)$.
We display backgrounds in the formulas so that this specialization is
visible. The lower equations are
\begin{align}
 \dd n_1&=0\quad\text{over }\Z,\\
 dn_2&=\omega_2\cup n_1+s_1\cup n_1^2\quad\text{over }\F,\\
 dn_3&=\mathcal O_4[n_1,n_2],\\
 \dd\widehat\nu_4&=\widehat O_5[n_1,n_2,n_3]\pmod1.
\end{align}
Binary occurrences of $n_1$ mean its parity; the integer itself is retained
where a lift carry or fractional phase requires it.

\section{Cochain arithmetic}
The lift $\lift z$ is taken after the entire expression $z$ is formed over
$\F$. Integer cup products use signed interval cuts. The differential on
phases is $\dd=d-2\lift{s_1}\cup$. For nonclosed cochains we use
\begin{equation}
 \SSq^j(u)=u\cup_{|u|-j}u+u\cup_{|u|-j+1}du,
 \qquad d\SSq^j(u)=\Sq^j(du).
\end{equation}
Negative cup indices give zero. The second term must not be omitted when
the lower equation makes $u$ nonclosed. Exact divisions precede reduction;
negative carries use floor rather than truncation toward zero.
The displayed obstruction and product formulas have denominators dividing
eight or sixteen. Their solved phase primitives are general exact rational
cochains modulo one, without this denominator restriction.

\section{The closed Majorana sector}
For a closed Majorana cochain $a$ of degree $p=1,2$ and a complex-fermion
cochain $c$ of degree $p+1$, set
\begin{equation}
 P(a)=\Sq^2a+\omega_2\cup a+s_1\cup\Sq^1a,
 \qquad dc=P(a).
\end{equation}
The total phase in the Gu--Wen (CA) coordinate is
\begin{equation}
 \widehat O_{p+3}(a,c)=\frac12\bigl[\SSq^2(c)+\omega_2\cup c\bigr]
 +\widehat O^\gamma_{p+3}(a).
\end{equation}
The full Majorana term is the current manuscript Cartan--Adem expression.
Its ordered words, Bockstein terms, and all antiunitary terms are evaluated
directly by the independent expression graph in \texttt{fspt/formulas.py}.
No microscopic phase table is queried.

For two states define
\begin{equation}
 m(a,b)=a\cup_{p-1}b+s_1\cup(a\cup_p b),\quad
 A=a+b,\quad C=c+c'+m(a,b).
\end{equation}
The phase product is
$\widehat\nu+\widehat\nu'+U_{p+2}(a,c;b,c')$.
The complete $U_3,U_4$ are evaluated as the six manuscript blocks:
Gu--Wen, Bockstein, Adem, extension, antiunitary, and mixed background.
They satisfy
\begin{equation}
 \dd U_{p+2}=\widehat O_{p+3}(A,C)
 -\widehat O_{p+3}(a,c)-\widehat O_{p+3}(b,c').
\end{equation}
This identity alone does not select a closed product correction.
The reference dictionary is retained when matching the physical operator
coordinate:
\begin{equation}
 \Gamma(c)=\tfrac12c\cup_{p+1}dc,\quad
 O_{\rm op}=O_{\rm CA}+\dd\Gamma,\quad
 U_{\rm op}=U_{\rm CA}+\Delta\Gamma.
\end{equation}
The source comparison also includes its specified closed unary phase $Z$.
It is not dropped because the obstruction is unchanged.

\section{The native integer-layer secondary source}
Put
\begin{equation}
 h_1=[\lfloor n_1/2\rfloor]_2,\qquad
 u=n_2+s_1\cup h_1,\qquad W=\omega_2+s_1^2.
\end{equation}
Then $du=W\cup n_1$, and the native short source is
\begin{align}
 \mathcal O_4[n_1,n_2]={}&\SSq^2(u)+\omega_2\cup u+s_1\cup\SSq^1(u)
 +\zeta_{2,1}(W,n_1)\nonumber\\
 &+(W\cup_1W+s_1\cup W)\cup h_1\nonumber\\
 &+\bigl[(W\cup_1W)\cup_1s_1+s_1\cup(s_1\cup_1W)\bigr]\cup n_1.
\end{align}
The Cartan term is the literal operation
$\zeta_{2,1}(W,a)=\smile_{123134}(W,W,a,a)$.
This fixes a cochain representative, not merely its cohomology class.
In particular the uniform higher-dimensional source is not substituted
without its accompanying coboundary dictionary.

\section{The complete integer-layer terminal phase}
Write $A=du$ and define the integer quotients
\begin{equation}
 K_3=\frac{\lift A-\lift W\cup n_1}{2},\qquad
 \beta^\circ u=\frac{d\lift u-\lift{du}}2.
\end{equation}
The second quotient is valid for nonclosed $u$; $d\lift u/2$ alone is not.
The terminal phase is assembled as
\begin{equation}
 \widehat O_5=\frac12[\SSq^2(n_3)+\omega_2\cup n_3]
 +\widehat O_5^{\gamma,\mathrm{ext}}(u)
 +\widehat O_5^{\gamma\psi}+\widehat O_5^\psi.
\end{equation}
Here the mixed term is
\begin{equation}
 \widehat O_5^{\gamma\psi}
 =\tfrac12[\SSq^2u+\omega_2\cup u+s_1\cup\SSq^1u]
 \cup_3\mathcal O_4^\psi[n_1].
\end{equation}
The pure integer term contains the fixed completion
$y_5=H_{\rm EZ}^*\mathcal R_6+AW^*y_5^{\rm tot}$, the quadratic $K_3$
terms, the whole lifted Cartan bracket, and the sixteenth-valued coupling
\begin{equation}
 \frac1{16}\bigl[\lift W\cup\lift W+
 \lift W\cup_1\dd\lift W\bigr]\cup n_1.
\end{equation}
All 55 fixed coefficients of $y_5^{\rm tot}$ and the complete individual
terms are those in the supplied normalized note. They are compilation
inputs, not unknown primitives to solve separately on each group.

The independent native evaluator compiles this entire formula into 7,239
scalar integer operations for $\omega_2=0$. The actual integer $n_1$,
both lower choices, transport signs, and carry bits are retained. The
compiled evaluator has no dependency on the old SptSet engine or on a
table of space-group classifications.
For the canonical torsion representative, the exact restriction is
\begin{equation}
 \widehat O_5(s_1,B,C)=\tfrac12\SSq^2(C)
          +\widehat O_5^{\gamma,\mathrm{ext}}(B)+\tfrac9{16}s_1^5.
\end{equation}
Indeed $h_1=0$, $\mathcal O_4^\psi[s_1]=0$, and the normalized pure phase
pulls back from $BC_2$ as $9s_1^5/16$. Every local sign simplex verifies
these restrictions against the unchanged supplied evaluator. A 935-instruction evaluator implements this case;
512 further local towers agree exactly with the generic expression.

\section{From cochains to classification and stacking}
The native free resolution is used for exact cocycles, primitives and
quotient coordinates. A primitive obtained there is corrected by the
comparison homotopy before it is used as a bar cochain. This correction
ensures that the displayed lower equations hold on the simplices where
the nonlinear formula is evaluated.

For the native pure-CF construction there is an additional comparison
between the two chosen Steenrod cochain representatives. Let $C$ be closed
and let $\nu_0$ be the bar image of the native phase seed, with the Majorana
and extension inputs zero. The exact reconstruction solves
\begin{equation}
 \rho_5=\bigl[2(O_{\rm bar}(0,C)-\dd\nu_0)\bigr]_2,
 \qquad dh_4=\rho_5,\qquad \nu_4=\nu_0+\tfrac12h_4\pmod1.
\end{equation}
Half-integrality and exact solvability are checked. The result satisfies
$\dd\nu_4=O_{\rm bar}(0,C)$, while $2\nu_4=2\nu_0$ modulo one.
Thus the comparison term cancels under doubling but is needed for the
actual lift. Agreement of native and bar operations in cohomology alone
does not identify their raw cochain vectors.

Changing the Majorana primitive changes the secondary source by the
appropriate lower operation. The $d_3$ test therefore includes all
relevant Majorana shifts; the $d_4$ test projects through the CF and
Majorana-secondary images. A single nonzero raw phase vector is not a
proof that the integer decoration is killed.

For the incoming gauge let $P=P(a_1)$. The state
$(0,P,O_4(a_1,0))$ has the required phase boundary. If $dc_2=P$, composing
it with the CF gauge differs from $O_4(a_1,c_2)$ by
\begin{equation}
 \tfrac12P\cup_2P=\dd\bigl[\tfrac12\SSq^1(c_2)\bigr].
\end{equation}
This exact witness aligns the group presentation with the chosen phase
coordinate.

The open-Majorana continuation also has a nontrivial zero-integer
coordinate dictionary. For $dB=0$ and $dC=\Sq^2B+s\Sq^1B$, we obtain
\begin{equation}
 O_5^{\rm pip}(0,B,C)-O_5^{\rm CA}(B,C)
 =\dd\left[D_0(s,B)+\tfrac14\lift{\Sq^2B}
                         +\tfrac12s\cup C\right].
\end{equation}
The quarter-valued local cochain $D_0$ solves all 32,768 closed-$B$
five-simplex equations. Another 128 complete legal towers verify the
identity. Thus the two formulas define the same secondary cohomology
class; their displayed phase cochains must not be identified pointwise.

\section{The torsion integer-layer square}
The supplied $U_3,U_4$ set the integer layer to zero. The collaborator's
calibrated cochain product is an additional mathematical input; its
universal coefficients are independently evaluated here. We align that
reference product with the normalized formulas before using its diagonal
on $n_1=s_1$, with $\omega_2=0$. Write $s=s_1$, $B=n_2$, $C=n_3$.
The CF coordinate dictionary is $C_{\rm ref}=C+\Lambda$, where
\begin{equation}
 \Lambda_{0123}=B_{013}+B_{012}B_{013}
                  +B_{012}B_{023}+B_{013}B_{023}.
\end{equation}
It satisfies $d\Lambda=Q_{\rm native}+Q_{\rm ref}$. The phase dictionary is
\begin{align}
 T_1(C)&=D_4(s,B)+\tfrac14s^4+\tfrac12
       [C\cup_2\Lambda+\Lambda\cup_3Q_{\rm native}],\\
 O_{\rm ref}(s,B,C+\Lambda)-O_{\rm native}(s,B,C)&=\dd T_1(C).
\end{align}
The fixed local cochain $D_4$ is tabulated on four sign edges and six
Majorana cone faces. Its 1,024 sixteenth-valued entries retain the initial
AW coordinate; the explicit $s^4/4$ converts this base to the corrected
supplied coordinate. Their sum solves all 32,768 legal local five-simplex
equations. The additional $C$ dependence follows
from the polarization identity for $\SSq^2$. No target-group answer enters
this universal cochain dictionary.

For the doubled output $n_1=2s,B=0$, one has $\Lambda_2=s^3$ and
\begin{equation}
 T_2(C)=-\tfrac3{16}s^4+\tfrac12 C\cup_2s^3.
\end{equation}
The restriction $B=0$ is part of this formula's stated domain.

To remove $2s$, let $t$ be the height zero-cochain on the cylinder and
$K$ its prism contraction to the bottom. The legal cylinder tower is
\begin{equation}
 n_I=2s+\dd t,\qquad B_I=s\cup t\cup dt,\qquad
 C_I=\pi^*C+KQ_{
 \rm native}(n_I,B_I).
\end{equation}
The endpoints are $(2s,0,C)$ and $(0,0,C+s^3)$. Its phase transport is
$g_4=\int_I O_5(n_I,B_I,C_I)$, satisfying
\begin{equation}
 \dd g_4=O_5(0,0,C+s^3)-O_5(2s,0,C).
\end{equation}
The resulting $g_4=13s^4/16+P(s,C)/2$ is an explicit binary polynomial in
four sign edges and four CF cone faces. Its complete coefficients are
recorded with the code; all 32,768 legal local five-simplex states verify
the displayed transport identity.

If the reference diagonal product has CF output $\beta$ and phase
correction $\Gamma_{\rm ref}$, its normalized lower-sector square is
\begin{equation}
 \left(0,\beta,\,
 2\widehat\nu+2T_1(C)+\Gamma_{\rm ref}
 -T_2(\beta+s^3)+g_4(\beta+s^3)\right).
\end{equation}
This retains the reference product's closed phase correction. The
extension order is subsequently obtained by reduction through the actual
lower-layer generators; no order-two lift is presumed.

\subsection{AW transport and phase convention}
Each integer edge in the AW middle block is transported from its own
initial vertex: $M_{ij}=(-1)^{s_{0i}}n_{ij}$. Correcting an earlier
compiler transcription changes the phase, at $\omega_2=0$, by
\begin{equation}
 \Delta O_5(012345)=\tfrac12 s_{01}s_{12}s_{23}s_{34}
 \operatorname{bit}_1\!\left((-1)^{s_{03}}n_{34}\right)
 (n_{45}\bmod2)\pmod1.
\end{equation}
This is zero for every even $n$ and for $s=0$. At $n=s$ it is
$s^5/2$, so the corrected flat phase is $\nu_{\rm new}=\nu_{\rm old}-s^4/4$.
The unary dictionary changes by $T_{1,\rm new}=T_{1,\rm old}+s^4/4$;
the canonical cylinder changes by $g_{\rm new}=g_{\rm old}+s^4/2$.
The reference $\Gamma_{\rm ref}$ is unchanged. Therefore the complete
square changes only by
\begin{equation}
 \nu_{\rm square,new}-\nu_{\rm square,old}
 =\tfrac12s^4=\dd\bigl[\tfrac12s\cup B\bigr],\qquad dB=s^3.
\end{equation}
This is an explicit phase gauge, proving equality of the marked relation
classes after the input lifts are rephased. Saved earlier cochains are
not current-coordinate witnesses. Direct comparison with the unchanged
supplied evaluator now includes 84 signed towers, negative integers up to
$2^{45}$, all sign simplices, and both original failing regression fixtures.

\section{Primitive free integer-layer generators}
The free rank alone does not fix the lattice of surviving integer
decorations. Put $H=H^1(G,\mathbb Z_s)$, and let $S\subset H$ be the
subgroup surviving all outgoing differentials. For $s=0$, every integral
character has the literal flat tower
\begin{equation}
 (n_1,n_2,n_3,\widehat\nu_4)=(n_1,0,0,0).
\end{equation}
In particular, substituting $s=n_2=n_3=0$ into the complete normalized
$O_5$ expression gives zero with $n_1$ left arbitrary. The native integral
cohomology basis therefore gives primitive free generators in this case.

For a signed character, the cocycle equation defines a homomorphism
\begin{equation}
 G\longrightarrow D_\infty=\mathbb Z\rtimes C_2,
 \qquad g\longmapsto(n_1(g),s(g)).
\end{equation}
Let $x$ denote the universal signed integer coordinate on $D_\infty$.
The class $2x$ has zero first source, so $n_2=0$ is allowed. Since
$D_\infty=C_2*C_2$, restriction to the two reflection groups is an
isomorphism in cohomological degrees at least two. On these subgroups,
$2x$ restricts respectively to $0$ and $2s$, and the normalized parity
source satisfies $Q(0,0)=Q(2s,0)=0$ pointwise. Thus a CF primitive exists.
Moreover,
\begin{equation}
 H^5(D_\infty,U(1)_s)
 =H^5(C_2,U(1)_s)\oplus H^5(C_2,U(1)_s)=0,
\end{equation}
because multiplication by two on $U(1)$ is surjective. The full phase
therefore has a primitive as well. Pullback proves $2H\subset S$.
An independent infinite-dihedral resolution constructs this complete
tower and verifies its defining equations with corrected bar primitives.

Consequently the primitive free lattice is determined by the joint
parity quotient $H/2H$, retaining the possible torsion shift of each
free character. Let $V$ be the projection of its surviving classes onto
the free quotient modulo two. Choose a reduced row-echelon basis
$v_i$ of $V$ and directly construct the flat tower of each row with its
surviving torsion shift. Adjoin $2e_j$ in the nonpivot columns, using the
universal doubled tower. These rows form an integral basis of the full
surviving lattice, of index $2^{r-\dim V}$ in the original rank-$r$
lattice. This construction requires no general integer-layer binary
product. Earlier files without these witnesses retain only an abstract
free splitting.

The accepted v09 archive \path{results/space_groups} contains complete
classification and stacking results for all 230 groups. All 44 groups with
nonzero free $p+ip$ rank export actual primitive lattice bases and complete
generator towers; all 32 surviving torsion $p+ip$ cases retain their full
upper witnesses. The archive manifest records the frozen source and
\texttt{normalized-pip-aw-edge-transport-v2} convention. This configuration
enables the closed-CF phase optimization and native mod-two contraction,
while retaining the integral bar comparison map. Its mathematical outputs
agree literally with the complete v07 baseline. Constructed witnesses and
their saved-certificate audits do not imply a separate full comparison-support
audit of every group; only cases explicitly recorded in the validation log
carry that additional numerical check.

\section{Examples distinguishing graded layers and extensions}
The new supplied space-group PDF lists four graded layers, with no stacking
extension column. Its 920 entries agree with the independent 230-group
graded calculation. A separate historical manuscript table contains four
stacking differences despite agreement of all four graded layers:
\begin{center}
\begin{tabular}{c|c|c}
 Space group & Historical extension & Independent extension\\\hline
 68 & $\mathbb Z_2^3\oplus\mathbb Z_4$ & $\mathbb Z_2^5$\\
 81 & $\mathbb Z\oplus\mathbb Z_2^4\oplus\mathbb Z_8^3$
    & $\mathbb Z\oplus\mathbb Z_2^2\oplus\mathbb Z_4\oplus\mathbb Z_8^3$\\
 82 & $\mathbb Z\oplus\mathbb Z_2^4\oplus\mathbb Z_8^2$
    & $\mathbb Z\oplus\mathbb Z_2^2\oplus\mathbb Z_4\oplus\mathbb Z_8^2$\\
 101 & $\mathbb Z_2\oplus\mathbb Z_4$ & $\mathbb Z_2^3$
\end{tabular}
\end{center}
For groups 68 and 101 there is one Majorana generator $B$ and a bosonic
kernel $\mathbb Z_2^k$, with $k=4,2$. The complete cochain reduction gives
$2B=0$, including nonzero incoming gauges and their phase primitives.
The raw square need not vanish pointwise. A nonzero $2B$ would give exactly
the historical alternative $\mathbb Z_4\oplus\mathbb Z_2^{k-1}$.
Because $2D=0$ for each boson $D$, changing the lift $B\mapsto B+D$
cannot account for the difference.

For groups 81 and 82 the measured marked relation is
\begin{equation}
 2P=C_1,\qquad 2C_1=0,
\end{equation}
with $C_1$ absent from all other lower relations. It gives a direct
$\mathbb Z_4$ summand. There is an independent geometric consistency
check: each actual affine group retracts onto its $C_4$ point group.
The quotient and an order-four affine section were verified exactly.
Consequently pullback and restriction force the $C_4$ classification to
be a direct summand of the space-group answer.

On $C_4$, write $s(a)=a\bmod2$, $\eta(a)=\lfloor a/2\rfloor$ and
$B(a,b)=s(a)\eta(b)$. Then $dB=s^3$, and the CF part of the square is
\begin{equation}
 K=B\cup_1B+B\cup_2dB+s\cup B,
 \qquad K(a,b,c)=(a\bmod2)\left\lfloor\frac{b+c}{4}\right\rfloor\bmod2.
\end{equation}
It evaluates to one on the mod-two bar cycle
$\sum_{i=0}^3[1|i|1]$, while $s^3$ evaluates to zero. This proves the
nonzero CF square independently of the upper-phase choice. The complete
finite calculation gives $\mathbb Z_4$, also listed in the supplied
$\omega_2=0$ point-group table. The historical alternatives
$\mathbb Z\oplus\mathbb Z_2^4\oplus\mathbb Z_8^k$ contain no direct
$\mathbb Z_4$ summand. These checks isolate an extension discrepancy;
they do not diagnose which operation produced the historical value.

\section{Verification and scope}
The independent Python and GAP evaluators are checked against exact
local lower towers and the supplied public formula evaluator. Separate
checks use actual bar primitives from disputed space groups. Local
arithmetic agreement, universal coefficient identities, and physical
normalization are recorded as different kinds of evidence.

The supplied normalized note proves its physical uniqueness statement
relative to a fixed lower tower and generator, with explicit finite-group
hypotheses. Using its cochains on a space group does not by itself extend
the stated theorem. Any additional topological identification required
for the crystalline application is to be stated separately from the
arithmetic calculation.
\end{document}
